% Labb 3, the final task, parts A and B, from labs/lab3/e_final_task.md. The wiring and the
% ready-made subroutines, headings of their own in the guide, are subsections; parts A and B are
% \labtask, with their Mål, Bakgrund, Specifikation, Uppgifter, Frågor and Extra as \task
% headings. The subroutines delay_s and delay_ms are the handout's, and are printed as the guide
% prints them.

\section{Slutuppgift: trafikljuset och övergångsstället}
\label{lab3:sec:e}
Slutuppgiften hör till kapitel~\ref{c11:ch}, och är den del av Labb~3 där programmet till sist
lämnar simulatorn och styr något på riktigt: tre lysdioder och en knapp på ett Arduino Uno-kort.
Del~A krävs för godkänt, och del~B för VG.

Läs \secref{c11:app:a} om lysdioder, knappar och kopplingen, och \secref{c11:app:b}, där ett
enklare ljus, ett vägarbetsljus, tas hela vägen från flödesplan till program. Hur programmet förs
över till kortet står i \secref{studio:5}.

\subsection{Kopplingen}
\label{lab3:sec:wiring}
Samma koppling används i båda delarna. Varje lysdiod har ett eget förkopplingsmotstånd, och knappen
kopplas mellan stift 2 och jord.

\begin{narrowtable}{@{}lcc@{}}
  \thead{Funktion} & \thead{Arduino-stift} & \thead{Port och bit} \\
  \midrule
  Röd lysdiod & 8 & PB0 \\
  Gul lysdiod & 9 & PB1 \\
  Grön lysdiod & 10 & PB2 \\
  Knapp, till jord & 2 & PD2 \\
  Jord & GND & -- \\
\end{narrowtable}

\bookfigure{l11_traffic_wiring}{Kopplingsschemat: tre lysdioder med förkopplingsmotstånd på stift
8, 9 och 10, och en knapp från stift 2 till jord.}{lab3:fig:wiring}

Lysdiodens långa ben, anoden, ska vara vänt mot motståndet och stiftet, och det korta, katoden, mot
jord. Koppla med USB-kabeln urkopplad, och kontrollera kopplingen mot tabellen innan kabeln sätts
i.

\subsection{De färdiga subrutinerna}
\label{lab3:sec:subroutines}
Båda delarna använder fördröjningar på hela sekunder. Ni behöver inte skriva dem själva: kopiera in
de två subrutinerna nedan sist i programmet, och konstanten \code{QUARTER_MS} överst. De förklaras
i \secref{c11:b4}, och \code{delay_ms} är samma subrutin som i del~22.

\begin{asmcode}
.equ QUARTER_MS = 250               ; 250 ms. Set it to 1 to run about 250 times
                                    ; faster in the simulator; set it back before
                                    ; flashing.
\end{asmcode}

\begin{asmcode}
;--------------------------------------------------------------------------------
; delay_s: Waits r24 seconds (1-63), as 4 x r24 calls of delay_ms with 250 ms each.
; A call takes 16 000 084 x r24 + 18 cycles: about 5 microseconds too long per second.
;--------------------------------------------------------------------------------
delay_s:
    push r24                        ; Save the registers this subroutine changes.
    push r25
    mov r25, r24
    lsl r25                         ; r25 = 2 x r24 ...
    lsl r25                         ; ... = 4 x r24: the number of quarter seconds.
    ldi r24, QUARTER_MS             ; A quarter of a second is 250 ms.
delay_s_loop:
    rcall delay_ms
    dec r25
    brne delay_s_loop
    pop r25                         ; Restore them, in the opposite order.
    pop r24
    ret

;--------------------------------------------------------------------------------
; delay_ms: Waits r24 milliseconds (1-255) at 16 MHz.
; A call takes 16 000 x r24 + 18 cycles, the rcall and the ret included.
;--------------------------------------------------------------------------------
delay_ms:
    push r24                        ; Save the registers this subroutine changes.
    push r26
    push r27
delay_ms_outer:
    ldi r26, low(3999)              ; X = r27:r26 = 3999.
    ldi r27, high(3999)
delay_ms_inner:
    sbiw r26, 1                     ; 4 cycles per turn, 3 the last time.
    brne delay_ms_inner
    dec r24                         ; One more millisecond done.
    brne delay_ms_outer
    pop r27                         ; Restore them, in the opposite order.
    pop r26
    pop r24
    ret
\end{asmcode}

Så används de: ladda \code{r24} med antalet sekunder och anropa \code{delay_s}. \code{ldi r24, 3}
följt av \code{rcall delay_s} väntar tre sekunder, och \code{r24} är fortfarande 3 efteråt.

\begin{aside}
\textbf{I simulatorn} tar varje simulerad sekund lång tid att köra. Sätt \code{QUARTER_MS} till 1
medan ni testar programmet där: då går varje \enquote{sekund} på ungefär 4~ms, och hela ljuscykeln
syns på några ögonblick. \textbf{Sätt tillbaka den till 250 innan programmet förs över till
kortet.}
\end{aside}

\labtask{Del A}{Trafikljuset}\phantomsection\label{lab3:a}
\task{Mål}
Styra ett trafikljus med tre lysdioder genom en fast ljuscykel, i simulatorn och på kortet.

\task{Bakgrund}
Vägarbetsljuset i \secref{c11:app:b} har två tillstånd; trafikljuset har fyra. Metoden är
densamma: skriv upp tillstånden och deras tider, rita flödesplanen, och skriv sedan programmet ett
tillstånd i taget.

\task{Specifikation}
Trafikljuset går igenom fyra tillstånd, i den här ordningen, och börjar sedan om:

\begin{narrowtable}{@{}lcccc@{}}
  \thead{Tillstånd} & \thead{Röd (PB0)} & \thead{Gul (PB1)} & \thead{Grön (PB2)} & \thead{Tid} \\
  \midrule
  Stopp & tänd & släckt & släckt & 3~s \\
  Klart att köra & tänd & tänd & släckt & 1~s \\
  Kör & släckt & släckt & tänd & 3~s \\
  Stopp strax & släckt & tänd & släckt & 1~s \\
\end{narrowtable}

\bookfigure{l11_traffic_timing}{Tidsdiagram för trafikljuset: en ljuscykel är 8~s.}%
{lab3:fig:timing}

\task{Uppgifter}
\begin{enumerate}
  \item \textbf{Flödesplan.} Rita flödesplanen för programmet, med ett symbolpar per tillstånd:
    tänd rätt lampor, vänta rätt tid. \kindtag[fillaccent]{Redovisa} den innan ni börjar skriva kod.
  \item \textbf{Bitmönstren.} Skriv upp vilket värde \code{PORTB} ska ha i vart och ett av de fyra
    tillstånden, binärt och hexadecimalt. Använd gärna \code{.equ RED = 0} och skrivsättet
    \code{(1 << RED)}, så att programmet säger vad det gör.
  \item \textbf{Programmet.} Skriv programmet utifrån programmallen, med initiering av \code{SP},
    \code{DDRB}, och de färdiga subrutinerna.
  \item \textbf{Simulera.} Sätt \code{QUARTER_MS} till 1 och kör i simulatorn. Sätt en brytpunkt på
    varje \code{out PORTB, r16} och kontrollera i I/O-fönstret att rätt lampor tänds, i rätt
    ordning. Kontrollera också med stoppuret att tiderna står i rätt förhållande till varandra: 3,
    1, 3, 1.
  \item \textbf{På kortet.} Sätt tillbaka \code{QUARTER_MS} till 250, bygg, och för över programmet
    till kortet enligt \secref{studio:5}.
  \item \textbf{Mät.} Ta tid på tio hela cykler med stoppuret i en mobiltelefon. \emph{Förutsäg}
    tiden först, med formeln för \code{delay_s}. Hur nära kommer ni? Vad begränsar noggrannheten i
    er mätning, och vad begränsar noggrannheten i programmet?
\end{enumerate}

\task{Frågor}
\begin{itemize}
  \item Vad händer om \code{delay_s} inte sparade \code{r24}? Hade programmet märkt det?
  \item Varför tänds röd och gul samtidigt i det andra tillståndet? Vad säger det till en bilist?
\end{itemize}

\redovisa{flödesplanen, simuleringen och trafikljuset på kortet.}

\labtask{Del B}{Övergångsstället}\phantomsection\label{lab3:b}
\task{Mål}
Låta en knapp styra när ljuscykeln körs: ett övergångsställe där bilarna har grönt tills en
fotgängare trycker på knappen.

\task{Bakgrund}
Att vänta på en knapp med \code{sbic} finns i \secref{c10:b5}, och varför en knapp studsar i
\secref{c11:a4}. Subrutiner och kursens regel om register finns i \secref{c10:app:c}.

\task{Specifikation}
\begin{itemize}
  \item Bilarna har \textbf{grönt}, så länge ingen trycker på knappen.
  \item När knappen trycks går ljuset igenom:

    \begin{narrowtable}{@{}lcccc@{}}
      \thead{Tillstånd} & \thead{Röd (PB0)} & \thead{Gul (PB1)} & \thead{Grön (PB2)} &
        \thead{Tid} \\
      \midrule
      Stopp strax & släckt & tänd & släckt & 1~s \\
      Stopp: fotgängarna går & tänd & släckt & släckt & 4~s \\
      Klart att köra & tänd & tänd & släckt & 1~s \\
    \end{narrowtable}
  \item Därefter blir det grönt igen, och programmet väntar på nästa tryck.
  \item Knappen är aktivt låg, med den inbyggda pull-upen: PD2 som ingång, med pull-up.
  \item Programmet ska använda minst en egen subrutin utöver de färdiga, och varje subrutin ska
    följa kursens regel: spara varje register den ändrar.
\end{itemize}

\task{Uppgifter}
\begin{enumerate}
  \item \textbf{Flödesplan.} Rita flödesplanen, med beslutet \enquote{knappen nedtryckt?} och de
    tre tillstånden. \kindtag[fillaccent]{Redovisa} den innan ni skriver kod.
  \item \textbf{En egen subrutin.} De tre tillstånden gör samma sak med olika värden: tänd vissa
    lampor, vänta en viss tid. Skriv en subrutin \code{show} som tänder lamporna i \code{r16} och
    väntar det antal sekunder som står i \code{r24}. Vilka register ändrar den? Måste något av dem
    sparas?
  \item \textbf{Programmet.} Skriv programmet. Initiera \code{SP}, \code{DDRB}, \code{DDRD} och
    pull-upen på PD2.
  \item \textbf{Simulera.} Sätt \code{QUARTER_MS} till 1. Kör programmet, stoppa det medan det
    väntar på knappen, och \enquote{tryck} genom att tömma rutan för bit~2 i \code{PIND}.
    \emph{Förutsäg} vilka lampor som ska tändas i vilken ordning, och kontrollera med
    brytpunkter.
  \item \textbf{På kortet.} Sätt tillbaka \code{QUARTER_MS} till 250 och för över programmet till
    kortet. Prova ett kort tryck, ett långt tryck, och ett tryck medan ljuset redan är rött.
\end{enumerate}

\task{Frågor}
\begin{itemize}
  \item Knappen studsar när den trycks. Varför ställer det inte till något här, när det gjorde det
    för knapptryckräknaren i \secref{c11:a4}?
  \item Vad händer om någon håller knappen nedtryckt hela tiden? Är det ett problem för bilarna?
  \item Ett tryck medan ljuset är rött gör ingenting. Varför inte? Hur skulle programmet behöva
    ändras för att komma ihåg ett sådant tryck?
\end{itemize}

\task{Extra}
Frivilliga utökningar, för den som vill vidare:
\begin{itemize}
  \item \textbf{Minsta tid på grönt.} Riktiga övergångsställen ger bilarna grönt i minst några
    sekunder innan ett nytt tryck får verkan. Låt bilarna alltid ha grönt i minst 5~s efter varje
    fotgängarfas.
  \item \textbf{Fotgängarnas ljus.} Lägg till en röd och en grön lysdiod för fotgängarna på PB3 och
    PB4 (stift 11 och 12). Fotgängarnas gröna ska lysa bara under de 4 sekunderna med rött för
    bilarna.
  \item \textbf{Vänta-lampan.} Tänd en lampa, till exempel PB5 (stift 13, lysdioden på kortet), när
    knappen har tryckts, och släck den när fotgängarna får grönt.
\end{itemize}

\redovisa{flödesplanen, simuleringen och övergångsstället på kortet.}
